% Book 3, Year 1, chapter 20 -- exhaustive search and backtracking.
% Provenance (never printed): sources/book3/BRIEFS.md, chapter 20.
\chapter{Exhaustive Search and Backtracking}\label{ch:b1:backtracking}

A sudoku that a puzzle page calls diabolical takes a person an evening. A program of forty lines solves it in a few
hundred steps, with no cleverness about sudoku at all: it fills an empty cell with a digit that breaks no rule, goes on,
and when it gets stuck it takes back its last choice and tries the next one. That method, \omterm{def:b1:backtracking:backtracking}{backtracking}, solves many
problems that have no better known algorithm: puzzles, schedules, \omterm{def:b1:propositional-logic:colouring}{colourings}, selections under a budget. Its cost can be
enormous; this chapter shows how to organise the search so that most of it is never done.

\begin{recall}
There are $2^n$ subsets of an $n$-element set, $n!$ orders of $n$ elements and $\binom nk$ subsets of $k$ elements.
Recursive functions and their termination were studied in \cref{ch:b1:induction-and-termination}.
\end{recall}

\section{Enumerating}\label{sec:b1:backtracking:enumerating}

\begin{definition}[Exhaustive search]\label{def:b1:backtracking:exhaustive}
An \emph{exhaustive search}\index{exhaustive search}, or \emph{brute force}\index{brute force}, solves a problem by
generating every candidate solution and testing each.
\end{definition}

The candidates are usually subsets, orders or selections, and the first task is to generate them all, each once. Subsets
of $n$ items correspond to the integers from 0 to $2^n - 1$ by their binary digits; they also come from a recursion that
decides the first item, in or out, then the rest.

\omcodes{code/bachelor-1/20-backtracking/listings/enumerate.py}{4-17}{Subsets by counting in binary, and by
deciding one item at a time}\label{lst:b1:backtracking:subsets}

\begin{omfigure}
\begin{tikzpicture}[n/.style={draw, rounded corners=2pt, font=\small\ttfamily, inner sep=2pt, minimum height=5mm},
  level distance=13mm, level 1/.style={sibling distance=60mm}, level 2/.style={sibling distance=30mm},
  level 3/.style={sibling distance=15mm}, e/.style={font=\scriptsize, midway, fill=white, inner sep=1pt}]
\node[n] {\{\}}
  child {node[n] {\{a\}}
    child {node[n] {\{a,b\}} child {node[n, fill=omDef!20] {abc}  edge from parent node[e] {+c}}
      child {node[n, fill=omDef!20] {ab} edge from parent node[e] {$-$c}} edge from parent node[e] {+b}}
    child {node[n] {\{a\}} child {node[n, fill=omDef!20] {ac} edge from parent node[e] {+c}}
      child {node[n, fill=omDef!20] {a} edge from parent node[e] {$-$c}} edge from parent node[e] {$-$b}}
    edge from parent node[e] {+a}}
  child {node[n] {\{\}}
    child {node[n] {\{b\}} child {node[n, fill=omDef!20] {bc} edge from parent node[e] {+c}}
      child {node[n, fill=omDef!20] {b} edge from parent node[e] {$-$c}} edge from parent node[e] {+b}}
    child {node[n] {\{\}} child {node[n, fill=omDef!20] {c} edge from parent node[e] {+c}}
      child {node[n, fill=omDef!20] {(none)} edge from parent node[e] {$-$c}} edge from parent node[e] {$-$b}}
    edge from parent node[e] {$-$a}};
\end{tikzpicture}
\par\smallskip{\small The include/exclude tree for $\{a, b, c\}$: each level decides one item; the eight leaves (shaded)
are the subsets, in the order the recursive generator yields them. The tree has $2^4 - 1 = 15$ nodes.}
\end{omfigure}

Orders are generated by putting each remaining element in the first free position, in turn, and ordering the rest;
swapping the candidate into place, and back afterwards, does it in the \omterm{def:b1:pointers-and-arrays:array}{array} itself. The lexicographic successor
computes the next order directly: find the longest decreasing suffix, exchange the element just before it with the
smallest larger element of the suffix, and reverse the suffix.

\omcodes{code/bachelor-1/20-backtracking/listings/enumerate.py}{20-43}{Orders by swapping, and the next order in
\omterm{def:b1:induction-and-termination:lex}{lexicographic order}}\label{lst:b1:backtracking:permutations}

\omoutput{code/bachelor-1/20-backtracking/outputs/enumerate.txt}{Subsets, orders and combinations of small
sets.}\label{out:b1:backtracking:enumerate}

Each generator spends a bounded amount of work per node of its tree, and its tree has a number of nodes proportional to
the number of outputs, times $n$ at most: the cost is proportional to the output, which is all one can ask. The output
itself is the problem: $2^{30}$ subsets is a billion, and $15!$ orders is more than a trillion.

\section{The tree of partial solutions}\label{sec:b1:backtracking:tree}

\begin{definition}[Partial solution, tree of partial solutions, backtracking]\label{def:b1:backtracking:backtracking}
A \emph{partial solution}\index{partial solution} is a sequence of choices that a complete solution may extend. The
\emph{tree of partial solutions}\index{tree of partial solutions} has the empty sequence as its root, and as children of
a node its extensions by one choice. \emph{Backtracking}\index{backtracking} explores this tree depth first, abandoning
a node, and its whole \omterm{def:b1:induction-and-termination:rooted-tree}{subtree}, as soon as the partial solution cannot be extended into a valid one.
\end{definition}

\begin{definition}[Constraint, pruning]\label{def:b1:backtracking:pruning}
The \emph{constraints}\index{constraint (of a search problem)} of a problem are the conditions a solution must satisfy.
\emph{Pruning}\index{pruning} is abandoning a node because its \omterm{def:b1:backtracking:backtracking}{partial solution} already breaks a constraint, or cannot
lead to a solution better than one already known.
\end{definition}

An \omterm{def:b1:backtracking:exhaustive}{exhaustive search} tests the \omterm{def:b1:backtracking:pruning}{constraints} at the leaves only; \omterm{def:b1:backtracking:backtracking}{backtracking} tests them as early as possible. The
difference is the number of nodes visited. The method itself is short, once the problem supplies its choices and its
tests:

\omcode{code/bachelor-1/20-backtracking/listings/backtrack.py}{4}{13}{The generic \omterm{def:b1:backtracking:backtracking}{backtracking}
method}\label{lst:b1:backtracking:generic}

Its demonstration builds the words of three letters over \code{a}, \code{b} with no two equal neighbours; the test on
the last two letters prunes every other branch at each level.

\omoutput{code/bachelor-1/20-backtracking/outputs/backtrack.txt}{The generic method on a small
problem.}\label{out:b1:backtracking:generic}

\section{Pruning: the $n$ queens}\label{sec:b1:backtracking:queens}

Place $n$ queens on an $n \times n$ chessboard so that no two attack each other: no two on a row, a column or a diagonal.
One queen goes on each row, so a placement is the list of the columns of rows $0, 1, \dots$; the columns must be
distinct, and two queens at $(r, c)$ and $(r', c')$ share a diagonal when $r - c = r' - c'$ or $r + c = r' + c'$. Sets of
the used columns and diagonals make each test $O(1)$.

\omcode{code/bachelor-1/20-backtracking/listings/queens.py}{7}{25}{The $n$ queens by \omterm{def:b1:backtracking:backtracking}{backtracking}, counting
nodes}\label{lst:b1:backtracking:queens}

\begin{omfigure}
\begin{tikzpicture}[n/.style={draw, rounded corners=2pt, font=\small\ttfamily, inner sep=2pt},
  d/.style={n, draw=omThm, text=omThm}, s/.style={n, fill=omDef!25, thick},
  level distance=10mm, level 1/.style={sibling distance=34mm}, level 2/.style={sibling distance=16mm}]
\node[n] {root}
  child {node[n] {0} child {node[d] {0 2 $\times$}} child {node[n] {0 3} child {node[d] {0 3 1 $\times$}}}}
  child {node[n] {1} child {node[n] {1 3} child {node[n] {1 3 0} child {node[s] {1 3 0 2}}}}}
  child {node[n] {2} child {node[n] {2 0} child {node[n] {2 0 3} child {node[s] {2 0 3 1}}}}}
  child {node[n] {3} child {node[n] {3 0} child {node[d] {3 0 2 $\times$}}} child {node[d] {3 1 $\times$}}};
\end{tikzpicture}
\par\smallskip{\small The \omterm{def:b1:backtracking:backtracking}{tree of partial solutions} for four queens, as visited: each node lists the columns of the
queens placed so far, row by row. Crossed nodes are dead ends: every cell of the next row is attacked. The two shaded leaves
are the solutions. Seventeen nodes, against 65 in the tree of all partial orders of the columns.}
\end{omfigure}

\omoutput{code/bachelor-1/20-backtracking/outputs/queens.txt}{Solutions and nodes visited, with \omterm{def:b1:backtracking:pruning}{pruning} and in the full
tree of placements with distinct columns.}\label{out:b1:backtracking:queens}

\begin{omfigure}
\begin{tikzpicture}
\begin{semilogyaxis}[width=10cm, height=6cm, xmin=3.5, xmax=10.5, xlabel={board size $n$}, ylabel={nodes visited},
  legend pos=north west, legend style={font=\small}, tick label style={font=\small}]
\addplot[thick, omThm, mark=*, mark size=1.2pt] table[x=n, y=full] {code/bachelor-1/20-backtracking/figures/nodes.dat};
\addlegendentry{full tree, $\sum_k n!/k!$}
\addplot[thick, omDef, mark=square*, mark size=1.2pt] table[x=n, y=pruned]
  {code/bachelor-1/20-backtracking/figures/nodes.dat}; \addlegendentry{with pruning}
\addplot[thick, black, dashed, mark=none] table[x=n, y=solutions]
  {code/bachelor-1/20-backtracking/figures/nodes.dat}; \addlegendentry{solutions}
\end{semilogyaxis}
\end{tikzpicture}
\par\smallskip{\small Nodes of the queens search on a logarithmic scale, from \cref{out:b1:backtracking:queens}. \omterm{def:b1:backtracking:pruning}{Pruning}
divides the work by 3.8 at $n = 4$ and by 278 at $n = 10$; both curves still grow exponentially.}
\end{omfigure}

\omterm{def:b1:backtracking:pruning}{Pruning} does not change the exponential growth, but it changes its rate, and that decides which sizes are reachable.
The tree without \omterm{def:b1:backtracking:pruning}{pruning} has $\sum_{k=0}^{n} n!/k! < e \cdot n!$ nodes (\cref{exo:b1:backtracking:9}).

\section{Sudoku and subset sum}\label{sec:b1:backtracking:sudoku}

Which choice to make first matters as much as when to test. In a sudoku, each empty cell can take the digits absent from
its row, its column and its $3 \times 3$ box. Choosing the cell with the fewest candidates makes dead ends appear early,
near the root, where abandoning a node discards a large \omterm{def:b1:induction-and-termination:rooted-tree}{subtree}; a cell with a single candidate costs no branching at
all.

\omcodes{code/bachelor-1/20-backtracking/listings/sudoku.py}{15-21,24-39}{Candidates of a cell, and the solver on the most
constrained cell}\label{lst:b1:backtracking:sudoku}

The grid was made by a program of this chapter that fills a grid at random and removes clues as long as the solution
stays unique (\code{make_grid.py}); a grid with a unique solution is what a puzzle page prints.

\omoutput{code/bachelor-1/20-backtracking/outputs/make_grid.txt}{The puzzle: 25 clues, the dots are the empty
cells.}\label{out:b1:backtracking:grid}

\omoutput{code/bachelor-1/20-backtracking/outputs/sudoku.txt}{The grid of \cref{out:b1:backtracking:grid}, solved with two rules for choosing the
next cell.}\label{out:b1:backtracking:sudoku}


The subset-sum problem asks for the subsets of given numbers with a given total: here, the reading lists of exactly 1000
pages from the library's books. With positive values, two tests prune the include/exclude tree: a partial sum already
above the target, and a partial sum that even all the undecided values cannot bring up to the target. Sorting the values
in decreasing order makes the first test fire early.

\omcode{code/bachelor-1/20-backtracking/listings/subset_sum.py}{15}{33}{Subset sum with two
\omterm{def:b1:backtracking:pruning}{prunings}}\label{lst:b1:backtracking:subset-sum}

\omoutput{code/bachelor-1/20-backtracking/outputs/subset_sum.txt}{Reading lists of 1000 pages from the first $n$
books.}\label{out:b1:backtracking:subset-sum}

Both tests are sound only because the values are positive (\cref{exo:b1:backtracking:6}): a \omterm{def:b1:backtracking:pruning}{pruning} rule must never
discard a \omterm{def:b1:induction-and-termination:rooted-tree}{subtree} that holds a solution.

\section{Branch and bound}\label{sec:b1:backtracking:bb}

The problems so far ask for all solutions, or one. An optimisation problem asks for the best: in the \emph{knapsack
problem}, items have weights and values, and the task is a set of items of total weight at most a capacity with the
largest total value.

\begin{definition}[Branch and bound, bounding function, continuous relaxation]\label{def:b1:backtracking:branch-and-bound}
\emph{Branch and bound}\index{branch and bound} explores the \omterm{def:b1:backtracking:backtracking}{tree of partial solutions} of an optimisation problem, keeps
the best complete solution found so far, and prunes a node when a \emph{bounding function}\index{bounding function}
shows that no completion of it can beat that solution. A \emph{continuous relaxation}\index{continuous relaxation} of a
problem allows fractions where the problem requires whole choices; its optimum is at least the problem's, and gives a
bound.
\end{definition}

For the knapsack, the relaxation allows a fraction of an item. Its optimum is easy: take the items in decreasing order of
value per unit of weight, and a fraction of the first one that does not fit. A node whose value so far plus the
relaxation's optimum on the remaining items and room does not exceed the best known value is pruned
(\cref{exo:b1:backtracking:11} proves that this is sound).

\omcodes{code/bachelor-1/20-backtracking/listings/knapsack_bb.py}{12-21,24-41}{The relaxation, and the search with and
without the bound}\label{lst:b1:backtracking:knapsack}

\omoutput{code/bachelor-1/20-backtracking/outputs/knapsack_bb.txt}{Nodes explored for $n$ random items and half their
total weight as capacity.}\label{out:b1:backtracking:knapsack}

\begin{omfigure}
\begin{tikzpicture}[n/.style={draw, rounded corners=2pt, font=\scriptsize, align=center, inner sep=2pt},
  pr/.style={n, draw=omThm, text=omThm}, lf/.style={n, fill=omDef!25, thick},
  e/.style={font=\scriptsize, midway, fill=white, inner sep=1pt},
  level distance=17mm, level 1/.style={sibling distance=52mm}, level 2/.style={sibling distance=30mm},
  level 3/.style={sibling distance=26mm}]
\node[n] {root\\value 0, room 10\\bound 76}
  child {node[n] {value 40, room 6\\bound 76}
    child {node[pr] {$7 > 6$:\\does not fit} edge from parent node[e] {take B}}
    child {node[n] {value 40, room 6\\bound 69}
      child {node[n] {value 65, room 1\\bound 69}
        child {node[lf] {leaf: value 65\\best 65} edge from parent node[e] {leave D}}
        edge from parent node[e] {take C}}
      child {node[pr] {bound 52\\$\leq 65$: pruned} edge from parent node[e] {leave C}}
      edge from parent node[e] {leave B}}
    edge from parent node[e] {take A}}
  child {node[pr] {value 0, room 10\\bound 57 $\leq 65$:\\pruned} edge from parent node[e] {leave A}};
\end{tikzpicture}
\par\smallskip{\small \omterm{def:b1:backtracking:branch-and-bound}{Branch and bound} on four items A, B, C, D of weights 4, 7, 5, 3 and values 40, 42, 25, 12
(decreasing value per unit of weight), capacity 10 (\code{figures/bb_tree.py}). The first \omterm{def:b1:induction-and-termination:rooted-tree}{leaf} reached, A and C, is the
optimum; its value 65 then prunes both remaining branches. Counted as \code{best} counts them (a pruned node is visited, an item that does not fit is not tried), plain \omterm{def:b1:backtracking:backtracking}{backtracking} visits 20 nodes here and \omterm{def:b1:backtracking:branch-and-bound}{branch and bound} 7.}
\end{omfigure}

Plain \omterm{def:b1:backtracking:backtracking}{backtracking} only prunes items that do not fit, and its count doubles with each item; with the bound, the count
barely grows here. These instances are friendly: on others the bound is weak and the search stays exponential. The
knapsack has another method, by dynamic programming, in \cref{ch:b1:dynamic-programming}.

\section{Exercises}\label{sec:b1:backtracking:exercises}

\begin{exercise}[$\star$]\label{exo:b1:backtracking:1}
In which order does \code{subsets_binary("abcd")} produce the subsets that contain \code{d}? And \code{subsets("abcd")}?
\end{exercise}

\begin{exercise}[$\star$]\label{exo:b1:backtracking:2}
Apply \code{next_permutation} to $[1, 3, 5, 4, 2]$, giving each step.
\end{exercise}

\begin{exercise}[$\star$]\label{exo:b1:backtracking:3}
Count the nodes of the queens search for $n = 4$ in the figure, by level. How many leaves are dead ends?
\end{exercise}

\begin{exercise}[$\star$]\label{exo:b1:backtracking:4}
For the subset sum of $\{7, 5, 4, 2\}$ (sorted) with target 9, which node is pruned first, and by which test?
\end{exercise}

\begin{exercise}[$\star\star$]\label{exo:b1:backtracking:5}
This generator forgets to swap back. Explain its output.
\omcode{code/bachelor-1/20-backtracking/listings/exo_perm_bug.py}{4}{10}{Orders without the undo}
\omoutput[2-2]{code/bachelor-1/20-backtracking/outputs/exo_perm_bug.txt}{On \code{abc}.}
\end{exercise}

\begin{exercise}[$\star\star$]\label{exo:b1:backtracking:6}
The \omterm{def:b1:backtracking:pruning}{pruning} ``the partial sum is already above the target'' is applied to values that can be negative. Explain the
output, and give a sound test for this case.
\omoutput[2-3]{code/bachelor-1/20-backtracking/outputs/exo_negative.txt}{Values $8, 5, -4, 3, -2$, target 6.}
\end{exercise}

\begin{exercise}[$\star\star$]\label{exo:b1:backtracking:7}
Write a generator of all the ways to write $n$ as a sum of distinct positive integers, with a \omterm{def:b1:backtracking:pruning}{pruning}, and check the
counts against another method.
\end{exercise}

\begin{exercise}[$\star\star$]\label{exo:b1:backtracking:8}
Write a word-search solver: is a word spelled in a grid of letters along a path of neighbouring cells (in eight
directions), each cell used at most once?
\end{exercise}

\begin{exercise}[$\star\star$]\label{exo:b1:backtracking:9}
Show that the tree of all partial orders of $n$ elements has $\sum_{k=0}^{n} n!/k!$ nodes, and that this is less than
$e \cdot n!$.
\end{exercise}

\begin{exercise}[$\star\star\star$]\label{exo:b1:backtracking:10}
Colour the vertices of a graph with three colours so that no edge joins two vertices of the same colour, by
\omterm{def:b1:backtracking:backtracking}{backtracking}, and test it against \omterm{def:b1:backtracking:exhaustive}{brute force}. Compare with the encoding of \omterm{def:b1:propositional-logic:colouring}{colouring} into a \omterm{def:b1:propositional-logic:sat}{satisfiability problem} of
\cref{ch:b1:propositional-logic}.
\end{exercise}

\begin{exercise}[$\star\star\star$]\label{exo:b1:backtracking:11}
Prove that the relaxation bound of the knapsack is at least the value of every completion of the \omterm{def:b1:backtracking:backtracking}{partial solution}, so
that \omterm{def:b1:backtracking:branch-and-bound}{branch and bound} still finds the optimum.
\end{exercise}

\begin{exercise}[$\star\star\star$]\label{exo:b1:backtracking:12}
For $n = 6$, the queens search visits 153 nodes. Using the left-right symmetry of the board, explain why it is $1 + 2(x
+ y + z)$ for three \omterm{def:b1:induction-and-termination:rooted-tree}{subtree} sizes, and find them by program.
\end{exercise}

\section{Problem: Eight queens and their mirrors}\label{sec:b1:backtracking:problem}

\begin{problem}[{Weekend problem --- Eight queens and their mirrors}]\label{pb:b1:backtracking:1}
The eight queens problem has 92 solutions. Many are rotations or reflections of others. A program lists the solutions by
\omterm{def:b1:backtracking:exhaustive}{brute force} over the $8!$ placements with distinct columns, applies to each the eight symmetries of the square, and
groups them.
\omoutput{code/bachelor-1/20-backtracking/outputs/queens_symmetry.txt}{The 92 solutions, up to the symmetries of the
board.}\label{out:b1:backtracking:symmetry}

\needspace{4\baselineskip}
\medskip\noindent\textbf{Part I --- The encoding.}
\begin{enumerate}
\item Why must a solution have exactly one queen on each row and each column?
\item Show that a solution is a permutation $p$ of $\{0, \dots, 7\}$ with all $r - p(r)$ distinct and all $r + p(r)$
distinct.
\item How many placements of 8 queens on 64 squares are there, and how many with one queen per row and column?
\item What are the ranges of $r - c$ and $r + c$, and how many diagonals of each direction are there?
\item Why does a set of used diagonals make the attack test constant-time?
\item Why does \omterm{def:b1:backtracking:exhaustive}{brute force} over $8!$ orders suffice here, though it would not for $n = 14$?
\end{enumerate}

\needspace{4\baselineskip}
\medskip\noindent\textbf{Part II --- Pruning.}
\begin{enumerate}[resume]
\item Read the nodes visited for $n = 8$ with and without \omterm{def:b1:backtracking:pruning}{pruning}, and their ratio.
\item Why do the solutions of $n = 6$ (4) number fewer than those of $n = 5$ (10)?
\item In what sense does the ratio of \cref{out:b1:backtracking:queens} grow?
\item Placing queens column by column instead of row by row, would the counts change?
\item Why is the first queen placed in all $n$ columns, while later queens have fewer candidates?
\item Estimate the nodes for $n = 12$ from the growth of the column.
\end{enumerate}

\needspace{4\baselineskip}
\medskip\noindent\textbf{Part III --- Symmetries of the square.}
\begin{enumerate}[resume]
\item List the eight symmetries of the square (rotations and reflections) as maps of cells $(r, c)$.
\item Show that each maps a solution to a solution.
\item The set of images of a solution is its \emph{orbit}. Show that two orbits are equal or disjoint.
\item Show that the size of an orbit is 8 divided by the number of symmetries that fix the solution.
\item Show that no solution of eight queens is fixed by a reflection.
\item Show that a solution fixed by the quarter-turn would have an orbit of at most 2 elements, and deduce from the output
which symmetries fix the symmetric solution.
\end{enumerate}

\needspace{4\baselineskip}
\medskip\noindent\textbf{Part IV --- Counting by program.}
\begin{enumerate}[resume]
\item What does the program check to call a permutation a solution?
\item How does it compute an image under a symmetry, and why does it sort the moved cells?
\item How are the orbits found, and why is each solution put in exactly one?
\item Check $11 \times 8 + 1 \times 4 = 92$.
\item Check on the printed symmetric solution that the half-turn fixes it.
\item State the named result, \emph{the eight-queens count}: 92 solutions, 12 up to symmetry, eleven orbits of size 8 and
one of size 4.
\end{enumerate}
\end{problem}
